\section{Communication}

Both faculty and students are 
overcommitted, overscheduled procrastinators, bless their hearts.
(The author is a faculty member and was once a student, so no judgment
is implied.)  We underestimated the importance of clear messaging
about CBTF procedures, especially since they were completely new for
most faculty and students.
We learned that unless you message both types of stakeholders early and often
about CBTF policies, and push pain back onto them if they don't
observe and enforce the policies, the resulting operational challenges
will generate enough extra work for CBTF to become a scaling bottleneck.

\subsection{Communication With Faculty}

\lesson{The CBTF is not a generalized proctor-outsourcing facility.}

For most students and faculty, the CBTF was an entirely new way of
doing exams.  The first courses using it were taught by faculty who
were familiar with (and indeed trying to emulate) existing successful
CBTFs, but as more courses began to use it, we had to systematically
disabuse faculty of many mistaken assumptions about how the CBTF
worked.
Our goal was to avoid having faculty fall into habits that would
interfere with scaling up a successful CBTF later.  The key points we
have had to reinforce multiple times include the following (they are
essentially all variations of the principles set forth in~\cite{fox-cbtf-not-proctoring}):

\begin{itemize}

\item You cannot ``block'' any  CBTF sessions for the exclusive use of
  your course, whether your course contains 5, 50, or 500 students.

\item Students cannot bring paper cheat sheets with them, even if you
  have allowed it.

\item Students cannot turn in their scratch paper to proctors with the
  goal of having it graded for partial credit alongside their
  computer-based exam.  (We initially allowed this, but as soon as
  more than 2 courses were using the CBTF, it became unwieldy.)

\item Your TAs cannot be the proctors, and therefore students cannot
  expect to ask questions during the exam.  (Our initial proctors were in
  fact TAs of one of the pioneering courses, but we needed to ensure
  faculty understood that that was incidental.)

\item Our proctors cannot come to your classrooms to proctor your
  exams: the CBTF is a single service, not a set of affordances that
  can be unbundled. 

\end{itemize}

\subsection{Communication With Students}

\lesson{It's the student's responsibility, not the CBTF staff's
  responsibility, to make and keep exam reservations in a timely way.}

CBTF exams were similarly new for most students.  They readily
understood the concept of making a reservation to take an exam, and
appreciated the flexibility.  Nonetheless,
the most important
messages for students, which continue to be an issue as we have scaled
up our CBTF efforts, are as follows:

\begin{itemize}

\item If you procrastinate making a reservation, and the result is
  that when you finally get around to it none of the available slots
  work for your schedule, that is your problem, not our problem.

\item If you make a reservation, fail to show up without cancelling
  and rescheduling, and are then unable to find any remaining slots
  within the exam window, that is your problem, not our problem.

\end{itemize}

In practice, we needed to define ``your problem'' to mean ``a
situation that causes friction for you and/or the course staff, rather
than for the CBTF itself.''  In particular, we should have been more
firm in stating to faculty that in these cases, the burden of
proctoring the student (in the classroom, remotely, or however), will
once again devolve onto them and their course staff.

This last point bears reemphasizing.  Many faculty are not great at
emphatically communicating important procedures to students; many
students are not great at retaining the message when they first hear
it.  By not being firm enough in our wording and repetitive enough in
nagging faculty to convey the message through multiple channels
(course home page, email announcement, posting to course question
boards, asking TAs to announce during office hours, and so on), we are
still suffering from higher rates than we'd like of students in the
above two scenarios who end up pushing a lot of extra work onto our
CBTF staff, thwarting scalability.  We are redesigning our
communication and doubling down on a structured and repetitive
``phased program'' of communication to students, to faculty, and to
students through faculty, to neutralize this operational challenge.
